\scanpage{136}
\mspage{67}
\[
(224)
\]
\diagram{\begin{tikzcd}[ampersand replacement=\&]
D_{d_0,\ldots,d_r}\arrow[r]\arrow[d]\&D_{d'_0,\ldots,d'_{r'}}\arrow[d]\\
D_{d_0,d'_0}\arrow[r]\&D_{d'_0}^{\deleted{*}}
\end{tikzcd}}
où les flèches verticales sont étales,
et la $2^{\mathrm{ième}}$ flèche horizontale est une
immersion loc.le de codim \deleted{$d_0-d'_0$. On a}
\[
\tag{225}d=d_0-d'_0
\]
[Dans le \unsure{typique} (auquel on se réduit)
où $d'_* = \{\uline{d_1}>\cdots>d_r\}$, le diagramme
(224) est cartésien -- \deleted{\unsure{comme on voit}}\inserted{plus génér.t}, il
reste \deleted{\unsure{tjs}} cartésien, \inserted{si $d'_0=d_1$}, et c'est trivial.]
\margin{\uline{NB} Dans le cas général, on a :
\diagram{\begin{tikzcd}[ampersand replacement=\&, column sep=small]
D_{d_*}\arrow[r,"{\text{fini étale}}"']\arrow[dd,dotted]\&D_{\widetilde d'_*}\arrow[r]\arrow[d]\&D_{d'_*}\arrow[d]\\
\&D_{d_0,d'_0}\arrow[dl,dotted]\arrow[r]\&D_{d'_0}\\
D_{d_0}\&\&
\arrow[ul,phantom,from=2-3,to=1-2,"{\text{cart.}}" description]
\end{tikzcd}}}
Considérons le morphisme induit
\[
(225)
\]
\diagram{\begin{tikzcd}[ampersand replacement=\&]D^*_{d_*}\arrow[r]\&D_{d'_*}\end{tikzcd}}
(Attention, on a mis un $*$ en exposant
à $D_{d_*}$ seulement, non à $D_{d'_*}$ !).

Nous munissons $D_{d'_*}$ du diviseur
$\Theta_{d'_*}$ image inverse de $\Theta_{d_0}$ sur $D_{d_0}$.
On désigne par $\Theta_{d'_*}^{(p)}$ le ``squelette
de dim $p$'' de $\Theta_{d'_*}^{p}$, et \deleted{\unsure{p}} on pose
\margin{135}
